\section{Factor-Group Computations and Simple Groups}

\exercise{15.1}{Removing an entire direct factor}
The projection $(a,b)\mapsto a$ from $\Z_2\times\Z_4$ onto $\Z_2$ has kernel $\{0\}\times\Z_4=\gen{(0,1)}$. The first isomorphism theorem therefore identifies the quotient with $\Z_2$.

\exercise{15.2}{Reducing the second coordinate modulo two}
The map $(a,b)\mapsto(a,b\bmod2)$ is a surjective homomorphism onto $\Z_2\times\Z_2$. Its kernel is $\{(0,0),(0,2)\}=\gen{(0,2)}$. Hence this is the quotient's classification.

\exercise{15.3}{A diagonal relation changes the quotient}
Define $\phi:\Z_2\times\Z_4\to\Z_4$ by $\phi(a,b)=2a+b\pmod4$. It is well defined because changing $a$ by 2 changes the result by 4, and it is onto since $\phi(0,1)=1$. Its kernel is exactly $(0,0),(1,2)$. Thus the quotient by $\gen{(1,2)}$ is $\Z_4$, rather than $\Z_2\times\Z_2$.

\exercise{15.4}{A cyclic quotient of order eight}
The surjective map $(a,b)\mapsto b-2a\pmod8$ from $\Z_4\times\Z_8$ to $\Z_8$ is well defined. Its kernel consists exactly of $(a,2a)$, the subgroup generated by $(1,2)$. Therefore the quotient is $\Z_8$. Its order is also $32/4=8$, checking the computation.

\exercise{15.5}{Eliminating one coordinate in a triple product}
Map $\Z_4\times\Z_4\times\Z_8$ onto $\Z_4\times\Z_8$ by
\[\phi(a,b,c)=(b-2a\pmod4,\ c-4a\pmod8).\]
Changing the representative for $a$ by 4 adds multiples of 4 and 8 to the outputs, so the formula is well defined. It is onto using $a=0$. Its kernel is $(a,2a,4a)$, precisely $\gen{(1,2,4)}$. Thus the quotient is $\Z_4\times\Z_8$, of order $128/4=32$.

\exercise{15.6}{An infinite coordinate projection}
The map $\Z^2\to\Z$, $(a,b)\mapsto a$, is onto with kernel $\gen{(0,1)}$. Hence the quotient is the infinite cyclic group $\Z$.

\exercise{15.7}{A primitive diagonal in the integer lattice}
Use $\phi(a,b)=b-2a$. This homomorphism onto $\Z$ has kernel the pairs $(a,2a)$, exactly $\gen{(1,2)}$. Thus the quotient is $\Z$, with no finite torsion factor.

\exercise{15.8}{A primitive diagonal in three dimensions}
The map $\phi:\Z^3\to\Z^2$ given by $(a,b,c)\mapsto(b-a,c-a)$ is onto and has kernel the triples $(a,a,a)$. Therefore the quotient by $\gen{(1,1,1)}$ is $\Z\times\Z$.

\exercise{15.9}{A finite relation on an infinite coordinate}
Reducing the first integer coordinate modulo 3 gives the quotient $\Z_3\times\Z\times\Z_4$. Since 3 and 4 are coprime, $\Z_3\times\Z_4\cong\Z_{12}$ by the Chinese remainder theorem. In invariant-factor form the answer is $\Z\times\Z_{12}$. The middle free coordinate remains infinite.

\exercise{15.10}{A fourth-multiple relation}
The relation only reduces the second integer coordinate modulo 4. The map $(a,b,c)\mapsto(a,b\bmod4,c)$ is onto $\Z\times\Z_4\times\Z_8$ with kernel $\gen{(0,4,0)}$. Thus that direct product is the classification, and $4\mid8$ puts the finite factors in invariant-factor order.

\exercise{15.11}{A nonprimitive diagonal}
Map $\Z^2$ onto $\Z\times\Z_2$ by $\phi(a,b)=(a-b,b\bmod2)$. Given $(u,\bar v)$, choose an integer representative $v$ and use $(u+v,v)$, proving surjectivity. The kernel requires $a=b$ and $b$ even, hence is $\gen{(2,2)}$. Therefore the quotient is $\Z\times\Z_2$. The factor of 2 in the diagonal produces the torsion.

\exercise{15.12}{A triple diagonal with torsion}
The map $(a,b,c)\mapsto(b-a,c-a,a\bmod3)$ is onto $\Z^2\times\Z_3$. Its kernel consists of $a=b=c$ with $3\mid a$, exactly $\gen{(3,3,3)}$. Thus the quotient is $\Z\times\Z\times\Z_3$.

\exercise{15.13}{Center and commutators of the square group}
Write $D_4=\gen{r,s}$ with $r^4=s^2=e$ and $srs=r^{-1}$. A rotation $r^k$ commutes with $s$ exactly when $r^k=r^{-k}$, meaning $k$ even; no reflection commutes with $r$, since that would give $r=r^{-1}$. Thus $Z(D_4)=\{e,r^2\}$.
The commutator $rsr^{-1}s^{-1}=r^2$ lies in $C$. On the other hand the quotient by $\gen{r^2}$ has $r^{-1}=r$, so its defining generators commute and it is abelian. The commutator subgroup is therefore contained in $\gen{r^2}$ as well, proving $C=\{e,r^2\}$.

\exercise{15.14}{A direct product with a symmetric group}
For any two groups, the center of their product is the product of their centers: commuting is checked in each coordinate. Likewise $[G\times H,G\times H]=[G,G]\times[H,H]$, since commutators are coordinatewise and commutators in either individual coordinate are obtained by setting the other coordinate to the identity. As $\Z_3$ is abelian, $Z(S_3)=\{e\}$, and $[S_3,S_3]=A_3$, the answers are
\[Z(\Z_3\times S_3)=\Z_3\times\{e\},\qquad C=\{0\}\times A_3\cong\Z_3.\]
For the last standard fact, $S_3/A_3$ is abelian, while the commutator of two different transpositions is a nonidentity three-cycle, generating $A_3$.

\exercise{15.15}{Two nonabelian direct factors}
Apply the coordinatewise center and commutator rules proved in Exercise 15.14. Since $Z(S_3)=\{e\}$, $Z(D_4)=\{e,r^2\}$, $[S_3,S_3]=A_3$, and $[D_4,D_4]=\gen{r^2}$,
\[Z(S_3\times D_4)=\{e\}\times\gen{r^2}\cong\Z_2,\qquad C=A_3\times\gen{r^2}\cong\Z_6.\]
The last isomorphism uses the coprime orders 3 and 2.

\exercise{15.16}{Every small subgroup of a sixteen-element group}
Let $G=\Z_4\times\Z_4$. Its complete list of subgroups of order at most 4, with quotients, is
\[
\begin{array}{c|c|c}
\text{subgroup}&\text{order}&G/H\\\hline
\{(0,0)\}&1&\Z_4\times\Z_4\\
\gen{(2,0)}&2&\Z_2\times\Z_4\\
\gen{(0,2)}&2&\Z_2\times\Z_4\\
\gen{(2,2)}&2&\Z_2\times\Z_4\\
\gen{(1,0)},\ \gen{(0,1)}&4&\Z_4\text{ (each)}\\
\gen{(1,1)},\ \gen{(1,2)}&4&\Z_4\text{ (each)}\\
\gen{(1,3)},\ \gen{(2,1)}&4&\Z_4\text{ (each)}\\
\{(0,0),(2,0),(0,2),(2,2)\}&4&\Z_2\times\Z_2
\end{array}
\]
Here each bracket denotes all multiples of its displayed generator, so the actual subgroup elements are explicit from that rule. To verify completeness, exactly three nonzero elements have order 2. The remaining twelve nonzero elements have order 4, with two generators per cyclic subgroup, giving exactly six cyclic subgroups as listed. A noncyclic order-four subgroup must consist of the identity and all three order-two elements, so there is exactly one more.
For an order-two subgroup, the quotient has order 8, exponent at most 4, and an element of order 4: use a coordinate not killed by that subgroup, or $(1,0)+H$ for $H=\gen{(2,2)}$. Its classification is therefore $\Z_2\times\Z_4$. Each cyclic order-four generator has an odd coordinate; elementary invertible changes of coordinates send it to $(1,0)$, so its quotient is $\Z_4$. The remaining subgroup is $2G$, and reducing both coordinates modulo 2 gives its quotient $\Z_2^2$.

\exercise{15.17}{The center must contain exactly the commuting elements}
Replace ``contains all'' by ``is the set of all.'' Precisely,
\[Z(G)=\{z\in G:zg=gz\text{ for every }g\in G\}.\]
Saying only that it contains these elements would allow extraneous elements and would not determine a subgroup uniquely.

\exercise{15.18}{The commutator subgroup requires generation}
The commutator subgroup is the subgroup \emph{generated} by all commutators, not merely their set. It is
\[C=\gen{a^{-1}b^{-1}ab:a,b\in G}.\]
This convention gives the same subgroup as $aba^{-1}b^{-1}$, by replacing $a,b$ with their inverses. Individual commutators need not themselves be closed under multiplication, so generation is essential.

\exercise{15.19}{Ten statements about factor and simple groups}
\begin{enumerate}[label=(\alph*)]
\item True. If $a$ generates $G$, then $aN$ generates every quotient $G/N$.
\item False. $\Z_2^2$ is not cyclic, but quotienting by one coordinate gives $\Z_2$.
\item False. $1/2+\Z$ has order 2.
\item True. $1/n+\Z$ has order $n$, including the identity when $n=1$.
\item False. An element killed by 4 must have representative $k/4$. Modulo integers there are only four such classes, and exactly two have order 4: $1/4+\Z,3/4+\Z$.
\item True. If every commutator is identity, $aba^{-1}b^{-1}=e$ gives $ab=ba$.
\item False. The inclusion is $C\subseteq H$, because all commutators become identity in an abelian quotient.
\item False. A simple cyclic group of prime order is abelian and has trivial commutator subgroup.
\item True. $C$ is normal; it is nontrivial in a nonabelian group, so simplicity forces $C=G$.
\item False. $A_5$ is simple and has order 60, which is not prime.
\end{enumerate}

\exercise{15.20}{Functions modulo constants}
Let $H=\{f:\R\to\R:f(0)=0\}$ under addition. Define $T:F\to H$ by $T(f)(x)=f(x)-f(0)$. This is an additive homomorphism and is onto because it fixes every function in $H$. Its kernel is exactly the constant functions $K$. Thus $F/K\cong H$. Choosing a basepoint turns each coset into its unique representative vanishing there.

\exercise{15.21}{Nonvanishing functions modulo constant factors}
Let $H^*=\{f:\R\to\R\setminus\{0\}:f(0)=1\}$ under multiplication. The map $T(f)(x)=f(x)/f(0)$ is well defined and multiplicative, and fixes every function in $H^*$, so is onto. Its kernel consists exactly of the nonzero constant functions. Therefore $F^*/K^*\cong H^*$.

\exercise{15.22}{No order-two additive coset modulo continuous functions}
No. If $f+K$ were killed by 2, then $2f$ would be continuous. Multiplying it by the constant $1/2$ would make $f$ continuous as well, so $f\in K$ and its coset is identity, not an element of order 2. The same argument shows the additive quotient has no nonidentity element of any finite order.

\exercise{15.23}{An order-two multiplicative coset}
Yes. Define $f(x)=1$ for $x\ge0$ and $f(x)=-1$ for $x<0$. It never vanishes but is discontinuous at 0, so $f\notin K^*$. Yet $f^2$ is the constant function 1 and lies in $K^*$. Thus $fK^*$ is nonidentity and its square is identity, giving order exactly 2. Unlike scalar multiplication, squaring loses the discontinuous sign information.

\exercise{15.24}{A translated full circle}
Since $z_0\in U$, multiplication by $z_0$ maps $U$ onto $U$: for any $w\in U$, $w=z_0(z_0^{-1}w)$ and $z_0^{-1}w\in U$. Hence $z_0U=U$, which is certainly a subgroup. Its quotient $U/z_0U=U/U$ is the trivial group.

\exercise{15.25}{The circle modulo a half-turn}
The homomorphism $U\to U$, $z\mapsto z^2$, is onto because every unit complex number has a unit square root. Its kernel is $\{1,-1\}=\gen{-1}$. Thus $U/\gen{-1}\cong U$ by the first isomorphism theorem.

\exercise{15.26}{The circle modulo the nth roots}
The map $U\to U$, $z\mapsto z^n$, is a surjective homomorphism. Its kernel is the cyclic subgroup of all $n$th roots of unity, generated by $\zeta_n$. Consequently $U/\gen{\zeta_n}\cong U$. When $n=1$ this is the identity-map case.

\exercise{15.27}{Real numbers modulo integers}
The map $\phi:\R\to U$ given by $\phi(t)=\cos(2\pi t)+i\sin(2\pi t)$ is a homomorphism from addition to multiplication, by the angle-addition formulas. It is onto, and $\phi(t)=1$ exactly when $t\in\Z$. Thus $\R/\Z\cong U$.

\exercise{15.28}{Torsion created by quotienting}
Take the torsion-free additive group $G=\Z$ and $H=2\Z$. In $G$, $na=0$ with $n>0$ forces $a=0$, so no nonzero element has finite order. But $G/H\cong\Z_2$, and all its elements have finite order. Alternatively $\Q/\Z$ gives an infinite quotient with the same property.

\exercise{15.29}{Isomorphic kernels with different quotients}
Let $G=\Z\times\Z_2$, $H=\Z\times\{0\}$, and $K=2\Z\times\{0\}$. They are normal because $G$ is abelian, and both are infinite cyclic, hence isomorphic. However $G/H\cong\Z_2$, while $G/K\cong\Z_2\times\Z_2$. Their different orders show the quotients are not isomorphic.

\exercise{15.30}{Centers of simple groups}
\begin{enumerate}[label=(\alph*)]
\item If a simple group is abelian, every element commutes with every other, so $Z(G)=G$. A nontrivial simple abelian group is cyclic of prime order.
\item If it is nonabelian, its normal center cannot be $G$. Simplicity leaves only $Z(G)=\{e\}$.
\end{enumerate}

\exercise{15.31}{Commutators of simple groups}
\begin{enumerate}[label=(\alph*)]
\item In an abelian group, every commutator is $e$, so $C=\{e\}$.
\item In a nonabelian group at least one commutator is nonidentity. Thus $C$ is a nontrivial normal subgroup, and simplicity forces $C=G$.
\end{enumerate}

\exercise{15.32}{Synopsis of cyclic quotients}
If $G=\gen a$, then every coset is $a^kN=(aN)^k$, so $aN$ generates $G/N$. This one observation proves the cyclic-quotient theorem.

\exercise{15.33}{Synopsis of maximal normal subgroups}
Under the quotient map $G\to G/M$, normal subgroups properly between $M$ and $G$ correspond exactly to nontrivial proper normal subgroups of $G/M$. Thus absence of the former, meaning $M$ is maximal normal, is equivalent to simplicity of $G/M$.

\exercise{15.34}{Index two prevents simplicity}
A subgroup $H$ of index 2 is normal: for $g\notin H$, both $gH$ and $Hg$ are the complement of $H$, while for $g\in H$ both equal $H$. The subgroup is proper because its index is 2 and nontrivial by hypothesis. It is therefore a nontrivial proper normal subgroup of $G$, so $G$ is not simple.

\exercise{15.35}{The image of a normal subgroup}
The image $\phi[N]$ is a subgroup of $\phi[G]$. For $\phi(g)\in\phi[G]$ and $\phi(n)\in\phi[N]$,
\[\phi(g)\phi(n)\phi(g)^{-1}=\phi(gng^{-1})\in\phi[N],\]
since $N$ is normal in $G$. Thus the image is normal in $\phi[G]$. It need not be normal in all of $G'$ when $\phi$ is not onto.

\exercise{15.36}{The inverse image of a normal subgroup}
Set $H=\phi^{-1}[N']$. It is a subgroup because identity, products and inverses map into the subgroup $N'$. For $g\in G,h\in H$,
\[\phi(ghg^{-1})=\phi(g)\phi(h)\phi(g)^{-1}\in N',\]
by normality of $N'$ in $G'$. Hence $ghg^{-1}\in H$, proving normality in $G$. No surjectivity assumption is needed.

\exercise{15.37}{A cyclic central quotient forces commutativity}
Prove the contrapositive. Suppose $G/Z(G)$ is generated by $aZ(G)$. Each $x\in G$ then has form $a^mz$ with $z\in Z(G)$, and similarly $y=a^nw$. Since the central factors commute with all elements,
\[xy=a^{m+n}zw=a^{n+m}wz=yx.\]
Thus $G$ is abelian. Therefore a nonabelian group cannot have a cyclic quotient by its center.

\exercise{15.38}{A nonabelian group of order a product of primes}
The center is a subgroup, so its order divides $pq$. If it were the whole group, $G$ would be abelian. If it were nontrivial and proper, its order would be $p$ or $q$ (also when $p=q$). The quotient would then have prime order and hence be cyclic. Exercise 15.37 would make $G$ abelian, again a contradiction. Thus $Z(G)=\{e\}$.

\exercise{15.39}{A complete proof of simplicity of the alternating group}
We compose permutations from right to left. All letters appearing distinctly in a displayed identity are assumed distinct. Put $c_i=(r,s,i)$ for fixed distinct $r,s$ and $i\notin\{r,s\}$. The source's description of ``$n$ special cycles'' means these $n-2$ actual three-cycles; repeated entries $i=r,s$ are not cycles.
\begin{enumerate}[label=(\alph*)]
\item A three-cycle is a product of two transpositions, for example $(a,b,c)=(a,c)(a,b)$, so is even and lies in $A_n$.
\item Every even permutation is a product of an even number of transpositions. Group them into adjacent pairs. Equal transpositions cancel. A pair sharing one letter is a three-cycle (or its inverse). For a disjoint pair, direct computation gives $(a,b)(c,d)=(a,c,b)(a,c,d)$. Thus every such pair is a product of three-cycles, and these generate $A_n$.
\item To show the special three-cycles suffice, compute
\begin{align*}
c_i^2&=(r,i,s),&c_jc_i^2&=(r,i,j),\\
c_j^2c_i&=(s,i,j),&c_i^2c_kc_j^2c_i&=(i,j,k).
\end{align*}
Their inverses give the opposite orientations. These cover three-cycles containing both $r,s$, only $r$, only $s$, or neither. Thus every three-cycle is a product of the special ones, and (b) proves generation. If $n=3$, only $c_i,c_i^{-1}$ are needed.
\item Suppose $N\trianglelefteq A_n$ contains a three-cycle, relabeled as $c_i$. Then $c_i^2\in N$. For any other $j\notin\{r,s,i\}$, the double transposition $h=(r,s)(i,j)$ is even, and conjugation gives $hc_i^2h^{-1}=(r,s,j)=c_j\in N$. Thus $N$ contains all special cycles (and $c_i$ itself). Part (c) gives $N=A_n$. For $n=3$ the contained three-cycle already generates $A_3$.
\item Take a nonidentity $\sigma\in N$ and use its disjoint-cycle decomposition. We check every case in the source. For any $t\in A_n$, normality gives $\sigma^{-1}t\sigma t^{-1}\in N$: the first factor $\sigma^{-1}$ and the conjugate $t\sigma t^{-1}$ both belong to $N$.
\step{Case I: $N$ contains a three-cycle as an element.} Part (d) applies immediately. Merely having a three-cycle among the disjoint factors of a larger permutation does not suffice; those cases are treated below.
\step{Case II: there is a cycle of length $\ell>3$.} Write $\sigma=\mu(a_1,a_2,\ldots,a_\ell)$, with $\mu$ disjoint from that cycle. Set $t=(a_1,a_2,a_3)$. Conjugation replaces the entries by their preimages under $\sigma$, so
\[\sigma^{-1}t\sigma t^{-1}=(a_\ell,a_1,a_2)(a_1,a_3,a_2)=(a_1,a_3,a_\ell).\]
This is a three-cycle in $N$, so (d) applies.
\step{Case III: at least two three-cycles occur.} Write $\sigma=\mu(a_4,a_5,a_6)(a_1,a_2,a_3)$, all factors disjoint, and choose $t=(a_1,a_2,a_4)$. Then
\[\sigma^{-1}t\sigma t^{-1}=(a_3,a_1,a_6)(a_1,a_4,a_2)=(a_1,a_4,a_2,a_6,a_3).\]
This five-cycle lies in $N$, reducing to Case II, hence to Case I.
\step{Case IV: precisely one three-cycle and some two-cycles occur.} Here $\sigma=\mu(a_1,a_2,a_3)$ with $\mu$ a disjoint product of transpositions. Since $\mu^2=e$ and disjoint factors commute, $\sigma^2=(a_1,a_3,a_2)\in N$, a three-cycle.
\step{Case V: all nontrivial cycles are transpositions.} Since $\sigma$ is even, there are an even number of them, at least two. Write $\sigma=\mu(a_3,a_4)(a_1,a_2)$ with $\mu$ disjoint. Taking $t=(a_1,a_2,a_3)$ gives
\[\sigma^{-1}t\sigma t^{-1}=(a_2,a_1,a_4)(a_1,a_3,a_2)=(a_1,a_3)(a_2,a_4)=\alpha\in N.\]
Because $n\ge5$, choose a letter $j$ different from these four and let $\beta=(a_1,a_3,j)\in A_n$. Another conjugation calculation gives
\[\beta^{-1}\alpha\beta\alpha=(j,a_1)(a_2,a_4)(a_1,a_3)(a_2,a_4)=(a_1,a_3,j)\in N.\]
This is again a three-cycle. These cases exhaust all possible nonidentity even permutations: a longer cycle gives II; otherwise the number of three-cycles is at least two, exactly one, or zero, giving III, IV, or V. Every nontrivial normal subgroup is therefore $A_n$, proving simplicity for $n\ge5$.
\end{enumerate}

\exercise{15.40}{The product of a subgroup and a normal subgroup}
The set $HN$ contains both subgroups. For $h_1n_1,h_2n_2\in HN$, normality gives
\[(h_1n_1)(h_2n_2)=h_1h_2(h_2^{-1}n_1h_2)n_2\in HN.\]
Also $(hn)^{-1}=n^{-1}h^{-1}=h^{-1}(hn^{-1}h^{-1})\in HN$. Together with the identity this proves $HN$ is a subgroup. Any subgroup containing both $H$ and $N$ contains every product $hn$, so contains $HN$. Thus $HN$ is the smallest such subgroup.

\exercise{15.41}{A product of two normal subgroups}
By Exercise 15.40, $NM$ is a subgroup. If $g\in G$ and $nm\in NM$, then $g(nm)g^{-1}=(gng^{-1})(gmg^{-1})$. The two factors lie in $N$ and $M$, respectively, by their normality, so the product lies in $NM$. Therefore $NM$ is normal in $G$.

\exercise{15.42}{Disjoint normal subgroups commute elementwise}
For $h\in H,k\in K$, let $c=hkh^{-1}k^{-1}$. It belongs to $K$, because $hkh^{-1}\in K$ and $k^{-1}\in K$. It also belongs to $H$, because $kh^{-1}k^{-1}\in H$ and $c=h(kh^{-1}k^{-1})$. The intersection hypothesis gives $c=e$, so $hk=kh$. Thus all pairs of elements from the two subgroups commute.
